Induction heads are the primary mechanism for in-context learning (ICL) in large language models, allowing them to copy patterns from the preceding context. While beneficial for few-shot learning, this mechanism may be overactive, leading to unintended "leakage" of pattern-consistent tokens even when the task explicitly requires stochastic or unrelated generation. In this work, we investigate whether induction heads drive such unintended pattern leakage. Using \gpttwo, we design a \taskname task where models are presented with repeating patterns followed by a request for a random word. Our results show that the model assigns disproportionately high probability (7.8\%) to pattern-completing tokens—approx. 3800$\times$ higher than uniform random probability. Through mechanistic intervention, we show that zero-ablating the top 5 induction heads reduces this leakage by over 95\%, with a single head (L5H1) contributing nearly 70\% of the effect. These findings provide causal evidence that the same circuitry enabling ICL is also a major driver of "over-patterning" and the "repetition curse" in language models, suggesting a fundamental trade-off between pattern sensitivity and stochastic control.
